% source: J. S. Milne, "Abelian Varieties" (course notes), v2.00, 2008, www.jmilne.org
% pdf_page: 0108
% doc_page: 102 (Chapter III, Section 5 — proof of Lemma 5.2, Proposition 5.3, Lemma 5.4)
% retrieved: 2026-07-16
% transcribed_by: claude-fable-5 (vision, rendered at 200dpi via pdftoppm)

% [running head: CHAPTER III. JACOBIAN VARIETIES, page 102]

\paragraph{Proof.}
(a) If $Q$ is not in the support of $D$, then
\[
  H^1(C, \mathcal{L}(D + Q))^{\vee} = \Gamma(C, \Omega^1(-D - Q))
\]
can be identified with the subspace of $\Gamma(C, \Omega^1(-D))$ of
differentials with a zero at $Q$. Clearly therefore we can take $U$ to be the
complement of the zero set of a basis of $H^1(C, \mathcal{L}(D))$ together
with a subset of the support of $D$.

(b) Let $D_0$ be the divisor zero on $C$. Then $h^1(D_0) = g$, and on
applying (a) repeatedly, we find that there is an open subset $U$ of $C^r$
such that $h^1(\sum P_i) = g - r$ for all $(P_1, \dots, P_r)$ in $U$. The
Riemann-Roch theorem now shows that
$h^0(\sum P_i) = r + (1 - g) + (g - r) = 1$ for all $(P_1, \dots, P_r)$ in
$U$. \hfill $\square$

In proving (5.1), we can assume that $k$ is algebraically closed. If $U'$ is
the image in $C^{(r)}$ of the set $U$ in (5.2b), then
$f^{(r)} \colon C^{(r)}(k) \to J(k)$ is injective on $U'(k)$, and so
$f^{(r)} \colon C^{(r)} \to W^r$ must either be birational or else purely
inseparable of degree $> 1$. The second possibility is excluded by part (b)
of the theorem, but before we can prove that we need another proposition.

\paragraph{Proposition 5.3.}
\textit{(a) For all $r \ge 1$, there are canonical isomorphisms}
\[
  \Gamma(C, \Omega^1) \xrightarrow{\;\simeq\;} \Gamma(C^r, \Omega^1)^{S_r}
  \xrightarrow{\;\simeq\;} \Gamma(C^{(r)}, \Omega^1).
\]
\textit{Let $\omega \in \Gamma(C, \Omega^1)$ correspond to
$\omega' \in \Gamma(C^{(r)}, \Omega^1)$; then for any effective divisor $D$
of degree $r$ on $C$, $(\omega) \ge D$ if and only if $\omega'$ has a zero at
$D$.}

\textit{(b) For all $r \ge 1$, the map
$f^{(r)*} \colon \Gamma(J, \Omega^1) \to \Gamma(C^{(r)}, \Omega^1)$ is an
isomorphism.}

A global 1-form on a product of projective varieties is a sum of global
1-forms on the factors. Therefore
$\Gamma(C^r, \Omega^1) = \bigoplus_i p_i^* \Gamma(C, \Omega^1)$, where the
$p_i$ are the projection maps onto the factors, and so it is clear that the
map $\omega \mapsto \sum p_i^* \omega$ identifies $\Gamma(C, \Omega^1)$ with
$\Gamma(C^r, \Omega^1)^{S_r}$. Because $\pi \colon C^r \to C^{(r)}$ is
separable,
$\pi^* \colon \Gamma(C^{(r)}, \Omega^1) \to \Gamma(C^r, \Omega^1)$ is
injective, and its image is obviously fixed by the action of $S_r$. The
composite map
\[
  \Gamma(J, \Omega^1) \to \Gamma(C^{(r)}, \Omega^1) \hookrightarrow
  \Gamma(C^r, \Omega^1)^{S_r} \simeq \Gamma(C, \Omega^1)
\]
sends $\omega$ to the element $\omega'$ of $\Gamma(C, \Omega^1)$ such that
$f^{r*} \omega = \sum p_i^* \omega'$. As $f^r = \sum f \circ p_i$, clearly
$\omega' = f^* \omega$, and so the composite map is $f^*$ which we know to be
an isomorphism (2.2). This proves that both maps in the above sequence are
isomorphisms. It also completes the proof of the proposition except for the
second part of (a), and for this we need a combinatorial lemma.

\paragraph{Lemma 5.4.}
\textit{Let $\sigma_1, \dots, \sigma_r$ be the elementary symmetric
polynomials in $X_1, \dots, X_r$, and let $\tau_j = \sum X_i^j \, dX_i$.
Then}
\[
  \sigma_m \tau_0 - \sigma_{m-1} \tau_1 + \dots + (-1)^m \tau_m
  = d\sigma_{m+1}, \text{ all } m \le r - 1.
\]

\paragraph{Proof.}
Let $\sigma_m(i)$ be the $m$th elementary symmetric polynomial in the
variables
\[
  X_1, \dots, X_{i-1}, X_{i+1}, \dots, X_r.
\]
Then
\[
  \sigma_{m-n} = \sigma_{m-n}(i) + X_i \sigma_{m-n-1}(i),
\]
% [proof continues on the next page]
